\section{Gamma and polygamma curvature}
Gamma, digamma, beta, and polygamma problems are governed by logarithmic derivatives.  The central move is to replace a difficult quotient or curvature assertion by an integral representation whose kernel has a controlled sign, or by an endpoint expansion that proves the proposed sign pattern cannot hold.

This section develops the curvature side of the theory.  The narrative is organized around reusable reductions: first the kernel or derivative identity, then the monotonicity or convexity consequence, and finally the application when the source problem is genuinely solved by that chain.  One-off endpoint checks are kept short in the main text and expanded only when they serve a later theorem or application.
\begin{theorem}
\label{res:simon-gamma-quotient-complement-cm}
For every \(0<\alpha<1\), the function \(1-\Gamma(x+\alpha)/(\Gamma(x)x^\alpha)\) is completely monotone on \((0,\infty)\).
\end{theorem}
\begin{proof}
For \(0<\alpha<1\), put
\[
F_\alpha(x)=\frac{\Gamma(x+\alpha)}{\Gamma(x)x^\alpha},\qquad x>0.
\]
Then \(F_\alpha\) is a Bernstein function. In fact, \(1-F_\alpha\) is completely monotone.

The stronger route asserting that \(F_\alpha\) is a complete Bernstein function is false.

Let \(\beta=1-\alpha\in(0,1)\). Since \(\Gamma(x+1)=x\Gamma(x)\), the beta integral gives
\[
F_\alpha(x)
=x^\beta\frac{\Gamma(x+\alpha)}{\Gamma(x+1)}
=\frac{x^\beta}{\Gamma(\beta)}
\int_0^\infty e^{-xu}(e^u-1)^{-\alpha}\,du.
\]
Also
\[
1=\frac{x^\beta}{\Gamma(\beta)}
\int_0^\infty e^{-xu}u^{-\alpha}\,du.
\]
Thus
\[
1-F_\alpha(x)
=\frac{x^\beta}{\Gamma(\beta)}
\int_0^\infty e^{-xu}\left(u^{-\alpha}-(e^u-1)^{-\alpha}\right)\,du.
\]

Define
\[
J_\alpha(t)=\int_0^t (t-u)^{\alpha-1}(e^u-1)^{-\alpha}\,du.
\]
By Laplace convolution,
\[
\mathcal L[J_\alpha](x)
=\Gamma(\alpha)x^{-\alpha}
\int_0^\infty e^{-xu}(e^u-1)^{-\alpha}\,du.
\]
Therefore
\[
F_\alpha(x)
=\frac{x}{\Gamma(\alpha)\Gamma(1-\alpha)}\mathcal L[J_\alpha](x).
\]

After the substitution \(u=tv\),
\[
J_\alpha(t)=
\int_0^1
(1-v)^{\alpha-1}
\left(\frac{t}{e^{tv}-1}\right)^\alpha\,dv.
\]
For fixed \(v\in(0,1]\), the function
\[
t\mapsto \frac{t}{e^{tv}-1}
\]
is decreasing because
\[
\frac{d}{dt}\frac{t}{e^{tv}-1}
=
\frac{e^{tv}-1-tv e^{tv}}{(e^{tv}-1)^2}\le0.
\]
The integrand is dominated by \((1-v)^{\alpha-1}v^{-\alpha}\), which is integrable on \((0,1)\). Hence \(J_\alpha\) is decreasing. Moreover,
\[
J_\alpha(0+)
=\int_0^1(1-v)^{\alpha-1}v^{-\alpha}\,dv
=\Gamma(\alpha)\Gamma(1-\alpha),
\]
and \(J_\alpha(t)\to0\) as \(t\to\infty\).

Let \(\mu_\alpha=-dJ_\alpha\). This is a positive finite measure with total mass
\[
\mu_\alpha([0,\infty))=\Gamma(\alpha)\Gamma(1-\alpha).
\]
Stieltjes integration by parts gives, for \(x>0\),
\[
\int_0^\infty e^{-xt}\,\mu_\alpha(dt)
=
\Gamma(\alpha)\Gamma(1-\alpha)-x\mathcal L[J_\alpha](x).
\]
Using the formula for \(F_\alpha\),
\[
1-F_\alpha(x)
=
\frac{1}{\Gamma(\alpha)\Gamma(1-\alpha)}
\int_0^\infty e^{-xt}\,\mu_\alpha(dt).
\]
Thus \(1-F_\alpha\) is completely monotone.

Consequently,
\[
F_\alpha'(x)
=
\frac{1}{\Gamma(\alpha)\Gamma(1-\alpha)}
\int_0^\infty t e^{-xt}\,\mu_\alpha(dt),
\]
so \(F_\alpha'\) is completely monotone. Since \(F_\alpha(x)>0\), \(F_\alpha\) is a Bernstein function.

For \(t>0\), differentiating the \(v\)-integral under the same endpoint domination gives \(\mu_\alpha(dt)=m_\alpha(t)\,dt\), where
\[
m_\alpha(t)
=
\alpha\int_0^1
(1-v)^{\alpha-1}
\left(\frac{t}{e^{tv}-1}\right)^\alpha
\left(
\frac{v e^{tv}}{e^{tv}-1}-\frac1t
\right)\,dv.
\]
The bracket is positive, because with \(a=tv>0\),
\[
\frac{a e^a}{e^a-1}>1.
\]
Therefore the Laplace-density route in the argument is also closed.

The case \(\alpha=1/2\) is the specialization
\[
F_{1/2}(x)=\frac{\Gamma(x+1/2)}{\Gamma(x)\sqrt{x}},
\]
and the same complement-CM representation proves that \(F_{1/2}\) is Bernstein.

For \(z=-r+i0\) with \(r\in(\alpha,1)\), take the upper-lip branch of \(z^{-\alpha}\). Then
\[
F_\alpha(-r+i0)
=
\frac{\Gamma(\alpha-r)}{\Gamma(-r)}r^{-\alpha}e^{-i\pi\alpha}.
\]
Here \(\Gamma(\alpha-r)<0\) and \(\Gamma(-r)<0\), so the real gamma ratio is positive. Hence
\[
\operatorname{Im} F_\alpha(-r+i0)<0.
\]
A complete Bernstein function has the Pick property and maps the upper half-plane into itself. Therefore \(F_\alpha\) is not a complete Bernstein function for \(0<\alpha<1\).
\end{proof}

\begin{proposition}[Laplace representation for \(V_q\)]
\label{res:baricz-vq-q-laplace-normal-form}
For every \(x>0\) and \(q>-1\), Baricz's function satisfies \(V_q(x)=\pi^{-1/2}\int_0^\infty s^{-1/2}e^{-x^2s}(1+s)^{-(q+1)}\,ds\).
\end{proposition}
\begin{proof}
Put \(u=t^2-x^2\). Then \(t=(u+x^2)^{1/2}\) and
\[
dt=\frac{du}{2(u+x^2)^{1/2}}.
\]
Since \(e^{x^2}e^{-t^2}=e^{-u}\), the factor \(2\) in the definition cancels the \(1/2\) from \(dt\), giving
\[
V_q(x)=\frac1{\Gamma(q+1)}
\int_0^\infty e^{-u}u^q(u+x^2)^{-1/2}\,du.
\]
This integral is finite for \(q>-1\): near \(0\) the integrand is \(O(u^q)\), and at infinity it is exponentially decaying.

For \(r>0\),
\[
r^{-1/2}=\frac1{\sqrt\pi}\int_0^\infty s^{-1/2}e^{-rs}\,ds.
\]
With \(r=u+x^2\), Tonelli's theorem applies because the integrand is nonnegative. Thus
\[
\begin{aligned}
V_q(x)
&=\frac1{\Gamma(q+1)\sqrt\pi}
\int_0^\infty\int_0^\infty
e^{-u}u^q s^{-1/2}e^{-(u+x^2)s}\,ds\,du\\
&=\frac1{\Gamma(q+1)\sqrt\pi}
\int_0^\infty s^{-1/2}e^{-x^2s}
\left(\int_0^\infty u^q e^{-(1+s)u}\,du\right)\,ds\\
&=\frac1{\sqrt\pi}\int_0^\infty
s^{-1/2}e^{-x^2s}(1+s)^{-(q+1)}\,ds.
\end{aligned}
\]
The last step uses
\[
\int_0^\infty u^q e^{-(1+s)u}\,du
=\frac{\Gamma(q+1)}{(1+s)^{q+1}}.
\]

This proves the claim.

Let \(a=q+1>0\). The normal form becomes
\[
V_{a-1}(x)=\frac1{\sqrt\pi}\int_0^\infty
s^{-1/2}e^{-x^2s}(1+s)^{-a}\,ds.
\]
For each integer \(n\ge0\), differentiation under the integral gives
\[
(-1)^n\frac{d^n}{da^n}V_{a-1}(x)
=\frac1{\sqrt\pi}\int_0^\infty
(\log(1+s))^n s^{-1/2}e^{-x^2s}(1+s)^{-a}\,ds\ge0.
\]
The differentiation is justified locally uniformly in \(a>0\): near \(s=0\), \((\log(1+s))^n s^{-1/2}=O(s^{n-1/2})\), and as \(s\to\infty\), the factor \(e^{-x^2s}\) dominates every logarithmic and polynomial factor.

Thus \(a\mapsto V_{a-1}(x)\) is completely monotone on \((0,\infty)\). Equivalently, after the change of variables \(y=\log(1+s)\), it is a positive Laplace transform in \(a\).

This proves the claim.

For \(a>0\), write
\[
F_x(a)=V_{a-1}(x)
=\int_{(0,\infty)} e^{-ay}\,d\nu_x(y),
\]
where \(\nu_x\) is the positive pushforward of
\[
\frac1{\sqrt\pi}s^{-1/2}e^{-x^2s}\,ds
\]
under \(y=\log(1+s)\). This measure is not concentrated at one point because it has positive density on \(s>0\), hence on \(y>0\).

For \(a_1\ne a_2\) and \(0<\lambda<1\), strict Holder inequality gives
\[
F_x(\lambda a_1+(1-\lambda)a_2)
<
F_x(a_1)^\lambda F_x(a_2)^{1-\lambda}.
\]
Strictness holds because \(e^{-a_1y}\) and \(e^{-a_2y}\) are not proportional \(\nu_x\)-almost everywhere unless \(\nu_x\) is supported at a single \(y\).

Thus \(a\mapsto F_x(a)\) is strictly log-convex on \((0,\infty)\). Since \(a=q+1\), the shift preserves strict log-convexity, so \(q\mapsto V_q(x)\) is strictly log-convex on \((-1,\infty)\) for every fixed \(x>0\).

This proves the claim.
\end{proof}

\begin{proposition}
\label{res:chen-choi-gurland-euler-sum-normal-form}
For \(F(x)=\Gamma(1/x)\Gamma(3/x)/\Gamma(2/x)^2\), \(\log F(x)=\log(4/3)+\sum_{m=1}^{\infty}A(mx)\), with locally uniform convergence of all derivative series on compact subintervals of \((0,\infty)\).
\end{proposition}
\begin{proof}
Let
\[
A(u)=2\log(u+2)-\log(u+1)-\log(u+3).
\]
Then \(A(u)>0\) for \(u>0\), since
\[
A(u)=\log\frac{(u+2)^2}{(u+1)(u+3)}
\]
and the numerator exceeds the denominator by \(1\).
Moreover
\[
A'(u)=\frac2{u+2}-\frac1{u+1}-\frac1{u+3}
=-\frac{2}{(u+1)(u+2)(u+3)}.
\]
Using
\[
\int_0^\infty e^{-ut}e^{-at}\,dt=\frac1{u+a},
\]
we get
\[
-A'(u)=\int_0^\infty e^{-ut}e^{-t}(1-e^{-t})^2\,dt.
\]
Since \(A(u)\to0\) as \(u\to\infty\), integration from \(u\) to infinity gives
\[
A(u)=\int_0^\infty e^{-ut}\frac{e^{-t}(1-e^{-t})^2}{t}\,dt.
\]
The density is positive, locally integrable at \(0\), and exponentially decaying at infinity. Hence \(A\) is strictly completely monotone.

Weierstrass' product gives
\[
\log\Gamma(z)=-\log z-\gamma z-
\sum_{m=1}^{\infty}\left(\log(1+z/m)-z/m\right).
\]
Put \(z=1/x\). In
\[
\log\Gamma(z)+\log\Gamma(3z)-2\log\Gamma(2z),
\]
the linear terms cancel. The logarithmic terms give \(\log(4/3)\). The summand at index \(m\) becomes
\[
\log\frac{(1+2z/m)^2}{(1+z/m)(1+3z/m)}
=\log\frac{(mx+2)^2}{(mx+1)(mx+3)}=A(mx).
\]
Thus
\[
\log F(x)=\log\frac43+\sum_{m=1}^{\infty}A(mx).
\]
For every compact \(x\in[\varepsilon,M]\) and every \(n\ge0\),
\[
\frac{d^n}{dx^n}A(mx)=m^nA^{(n)}(mx)=O_{n,\varepsilon}(m^{-2}),
\]
because the three-point second difference in \(A^{(n)}\) is \(O((mx)^{-n-2})\). Therefore the series may be differentiated termwise locally uniformly.

For \(n=0\), \(\log F(x)>0\) follows from \(\log(4/3)>0\) and \(A(mx)>0\). For \(n\ge1\),
\[
(-1)^n(\log F)^{(n)}(x)=
\sum_{m=1}^{\infty}m^n(-1)^nA^{(n)}(mx)>0.
\]
Therefore \(\log F\) is strictly completely monotone on \((0,\infty)\), proving Chen-Choi Conjecture 1.
\end{proof}

\begin{proposition}[Bernstein range of the W-symbol]
\label{res:wakrim-w-symbol-exact-bf-range}
For \(0<\alpha<1\) and \(\beta\ge0\), the Wakrim W-symbol \(\Phi_{\alpha,\beta}(s)=s^\alpha(1+(1-\alpha)s^{\alpha-1})^{-\beta}\) is a Bernstein function on \((0,\infty)\) if and only if \(0\le\beta\le1\).
\end{proposition}
\begin{proof}
Set \(a=1-\alpha\in(0,1)\). Then
\[
\Phi_{\alpha,\beta}(s)
=s^{1-a}\left(1+a s^{-a}\right)^{-\beta}
=\frac{s^{1-a+a\beta}}{(s^a+a)^\beta}.
\]
Write
\[
p=1-a+a\beta=\alpha+(1-\alpha)\beta.
\]
Differentiating gives
\[
\Phi_{\alpha,\beta}'(s)
=s^{p-1}(s^a+a)^{-\beta-1}
\left[p(s^a+a)-\beta a s^a\right].
\]
Since \(p-\beta a=1-a=\alpha\),
\[
\Phi_{\alpha,\beta}'(s)
=s^{-a(1-\beta)}(s^a+a)^{-\beta-1}
\left[\alpha s^a+a(\alpha+a\beta)\right].
\]
With \(y=s^a\), this becomes
\[
\Phi_{\alpha,\beta}'(s)=H(s^a),
\]
where
\[
H(y)=
y^{\beta-1}(y+a)^{-\beta}
\left(\alpha+\frac{a^2\beta}{y+a}\right).
\]

For \(0<\beta\le1\), each factor in \(H\) is completely monotone on \((0,\infty)\): \(y^{\beta-1}=y^{-(1-\beta)}\), \((y+a)^{-\beta}\), and the positive sum \(\alpha+a^2\beta/(y+a)\). Therefore \(H\) is completely monotone by product and positive-sum closure.

The map \(s\mapsto s^a\) is Bernstein for \(0<a<1\). The standard composition theorem says that if \(f\) is completely monotone and \(g\) is Bernstein, then \(f\circ g\) is completely monotone. Hence \(\Phi_{\alpha,\beta}'\) is completely monotone, and \(\Phi_{\alpha,\beta}\) is Bernstein for \(0<\beta\le1\). The endpoint \(\beta=0\) is \(\Phi_{\alpha,0}(s)=s^\alpha\), also Bernstein.

Combining this local positive range with Wakrim's source proof that \(\Phi_{\alpha,\beta}\) is not Bernstein for \(\beta>1\), the exact range is
\[
\Phi_{\alpha,\beta}\in BF
\quad\Longleftrightarrow\quad
0\le\beta\le1.
\]
\end{proof}

\begin{proposition}
\label{res:chiu-yin-cm-derivative-reflection-lemma}
If \(g\) is completely monotone on \((0,\infty)\), then \(-g'\) is completely monotone wherever the derivative exists.
\end{proposition}
\begin{proof}
Chiu--Yin Remark 3.5 conjectures the converse of Lemma 3.1: in the Sparre Andersen model, complete monotonicity of the ascending ladder-height density should imply complete monotonicity of the claim-size density. This pass proves only the classical Cramer--Lundberg/equilibrium-ladder slice.

\emph{Setup.}
In the classical Cramer--Lundberg/equilibrium formulation, let \(H\) be the claim-size distribution, let
\[
\overline H(x)=1-H(x),
\]
and assume \(H\) is absolutely continuous on \((0,\infty)\). Let
\[
m=\int_0^\infty \overline H(u)\,du
\]
be the finite mean, with \(0<m<\infty\). The equilibrium excess cdf and density are
\[
F_e(x)=\frac1m\int_0^x\overline H(u)\,du,
\qquad
f_e(x)=F_e'(x)=\frac{\overline H(x)}{m}.
\]
In the classical ladder normalization under discussion, the ascending ladder-height density is this equilibrium density:
\[
f_+(x)=f_e(x)=\frac{\overline H(x)}{m}.
\]

If \(g\) is completely monotone on \((0,\infty)\), then for every \(n\ge0\),
\[
(-1)^n g^{(n)}(x)\ge0.
\]
Therefore
\[
(-1)^n(-g')^{(n)}(x)=(-1)^{n+1}g^{(n+1)}(x)\ge0,
\]
so \(-g'\) is completely monotone.

Assume \(f_+\) is completely monotone. Then \(\overline H=m f_+\) is completely monotone. Define the smooth density representative
\[
h_*(x)=-\overline H'(x)=-m f_+'(x).
\]
By the reflection lemma, \(-f_+'\) is completely monotone; multiplying by \(m>0\) preserves complete monotonicity. Hence \(h_*\) is completely monotone.

Since densities are a.e. representatives, the conclusion is that \(H\) admits a completely monotone density version \(h_*\). Thus the converse of Chiu--Yin Lemma 3.1 is true in the classical equilibrium-ladder slice in this density-version sense.

This does not solve the full Sparre Andersen converse in Remark 3.5. The general case includes the weak descending-ladder renewal factor in the representation of \(F^+\), so the original claim density is not recovered by a single derivative of the ladder-height density.
\end{proof}

The next result applies the preceding method to the fixed source problem \textbf{APP-0021}.
\begin{corollary}[APP-0021: Szabo's cutoff problem has \(\alpha\_0=1\)]
\label{res:dagum-c-endpoints-and-upper-gate}
For the Dagum threshold \(c(\beta)=\inf\{\alpha\ge0:x^{-\alpha}(1+x^\beta)^{-1}\in CM\}\), one has \(c(1)=0\), \(c(2)=1\), and \(c(\beta)\le\beta/2\) for \(1<\beta\le2\).

This resolves the fixed application label \textbf{APP-0021}: Determine the exact cutoff \(\alpha\_0\) for \(y^\alpha H\_d(y)\).
\emph{Source reference.} \cite{app-app-0021-source}.
\end{corollary}
\begin{proof}
The source theorem proves complete monotonicity for \(\alpha\le1\).  Conversely,
as \(y\downarrow0\),
\[
  H_d(y)={1-d\over y}+O(1).
\]
Hence, for \(\alpha>1\),
\[
  y^\alpha H_d(y)=(1-d)y^{\alpha-1}+O(y^\alpha),
\]
whose derivative is positive near zero.  A completely monotone function must be
non-increasing, so no \(\alpha>1\) is admissible.
\end{proof}

\begin{proposition}[Szabo psi difference exact alpha0 equals one]
\label{res:szabo-psi-difference-alpha0-exact-one}
For \(0<d<1\), the sharp cutoff for complete monotonicity of \(y^\alpha[\psi(y+d)-\psi(y)-d/y]\) is \(\alpha_0=1\).
\end{proposition}
\begin{proof}
Let \(0<d<1\), \(y>0\), and
\[
H_d(y)=\psi(y+d)-\psi(y)-\frac{d}{y}.
\]
Then
\[
y^\alpha H_d(y)
\]
is strictly completely monotone on \((0,\infty)\) if \(\alpha\le 1\), and is not completely monotone if \(\alpha>1\). Hence the sharp cutoff in Szabo Open Problem 1.5 is
\[
\alpha_0=1.
\]

As \(y\downarrow0\), the digamma Laurent expansion gives
\[
\psi(y)=-\frac1y-\gamma+O(y),
\qquad
\psi(y+d)=\psi(d)+O(y).
\]
Therefore
\[
H_d(y)=\psi(y+d)-\psi(y)-\frac{d}{y}
=\frac{1-d}{y}+O(1).
\]
Since \(0<d<1\), the leading coefficient is positive. Hence
\[
y^\alpha H_d(y)=(1-d)y^{\alpha-1}+O(y^\alpha).
\]
If \(\alpha>1\), then
\[
\frac{d}{dy}\left(y^\alpha H_d(y)\right)
=(1-d)(\alpha-1)y^{\alpha-2}+O(y^{\alpha-1})>0
\]
for all sufficiently small \(y>0\). A completely monotone function must be nonincreasing, so no \(\alpha>1\) can be completely monotone.

Combining this with the source-imported strict complete monotonicity for \(\alpha\le1\) proves \(\alpha_0=1\).

Szabo's original function is recovered by \(y=x+a\) and \(d=b-a\), where \(a\ge0\), \(b>0\), and \(0<b-a<1\). Translation from \((-a,\infty)\) to \((0,\infty)\) preserves complete-monotonicity sign tests.
\end{proof}

\begin{proposition}
\label{res:bs-c-half-catalan-hausdorff-representation}
For the Bondesson--Steutel distribution at \(c=1/2\), \(P_n(1/2)=C_n/(2\cdot4^n)=\frac{1}{\pi}\int_0^1 x^{n-1/2}(1-x)^{1/2}\,dx\), where \(C_n\) is the \(n\)-th Catalan number.
\end{proposition}
\begin{proof}
For \(c=1/2\), the source pgf is
\[
P(z)=\frac{1-\sqrt{1-z}}{z}.
\]
Using the Catalan generating function
\[
\sum_{n\ge0}C_n x^n=\frac{1-\sqrt{1-4x}}{2x},
\]
with \(x=z/4\), we get
\[
P(z)=\frac12\sum_{n\ge0}C_n\left(\frac z4\right)^n.
\]
Hence
\[
P_n(1/2)=\frac{C_n}{2\cdot4^n}
=\frac{1}{2\cdot4^n(n+1)}\binom{2n}{n}.
\]

Now
\[
\frac1\pi\int_0^1 x^{n-1/2}(1-x)^{1/2}\,dx
=\frac1\pi B(n+1/2,3/2).
\]
Since \(\Gamma(3/2)=\sqrt\pi/2\) and
\[
\Gamma(n+1/2)=\frac{(2n)!\sqrt\pi}{4^n n!},
\]
this beta integral equals
\[
\frac{(2n)!}{2\cdot4^n n!(n+1)!}
=\frac{C_n}{2\cdot4^n}
=P_n(1/2).
\]
Therefore
\[
P_n(1/2)=\frac1\pi\int_0^1 x^{n-1/2}(1-x)^{1/2}\,dx.
\]
This is a Hausdorff moment representation on \([0,1]\), so \((P_n(1/2))\) is completely monotone as a sequence.

For \(c=1/2\), the canonical sequence is
\[
r_n=(n+1/2)P_n(1/2).
\]
Using the Hausdorff representation,
\[
r_n=\frac{n+1/2}{\pi}\int_0^1 x^{n-1/2}(1-x)^{1/2}\,dx.
\]
Because
\[
\frac{d}{dx}x^{n+1/2}=(n+1/2)x^{n-1/2},
\]
integration by parts gives
\[
r_n
=\frac1\pi\int_0^1 (1-x)^{1/2}\,d(x^{n+1/2})
=\frac{1}{2\pi}\int_0^1 x^{n+1/2}(1-x)^{-1/2}\,dx.
\]
The boundary terms vanish at \(0\) and \(1\). Thus
\[
r_n=\int_0^1 x^n\,d\nu(x),
\qquad
d\nu(x)=\frac{1}{2\pi}x^{1/2}(1-x)^{-1/2}\,dx.
\]
This is a positive finite measure on \([0,1]\). Hence \((r_n)\) is completely monotone.
\end{proof}

\begin{proposition}
\label{res:simon-leroy-moment-property-normal-form}
For \(\alpha,\beta,\gamma>0\), if Simon's moment property \(M_{\alpha,\beta,\gamma}\) holds for a positive random variable \(X\), then \((\Gamma(\beta))^\gamma F^{(\gamma)}_{\alpha,\beta}(z)={\bf E}[e^{zX}]\) for every \(z\in\mathbb C\).
\end{proposition}
\begin{proof}
\emph{Setup.}
Let
\[
F^{(\gamma)}_{\alpha,\beta}(z)
=\sum_{n\ge0}\frac{z^n}{\Gamma(\beta+\alpha n)^\gamma},
\qquad \alpha,\beta,\gamma>0.
\]

Assume Simon's moment property \(M_{\alpha,\beta,\gamma}\): there is a positive random variable \(X\) with
\[
{\bf E}[X^s]
=\Gamma(1+s)\left(\frac{\Gamma(\beta)}{\Gamma(\beta+\alpha s)}\right)^\gamma,
\qquad s>0.
\]

For \(n\ge0\), this gives
\[
{\bf E}[X^n]
=n!\left(\frac{\Gamma(\beta)}{\Gamma(\beta+\alpha n)}\right)^\gamma.
\]

Set
\[
a_n=\left(\frac{\Gamma(\beta)}{\Gamma(\beta+\alpha n)}\right)^\gamma.
\]
By the standard Gamma-ratio asymptotic,
\[
\frac{a_{n+1}}{a_n}
=\left(\frac{\Gamma(\beta+\alpha n)}{\Gamma(\beta+\alpha(n+1))}\right)^\gamma
\sim (\alpha n)^{-\alpha\gamma}\to0.
\]
Hence \(\sum_{n\ge0}a_n z^n\) has infinite radius of convergence.

For \(t\ge0\), monotone convergence gives
\[
{\bf E}[e^{tX}]
=\sum_{n\ge0}\frac{t^n{\bf E}[X^n]}{n!}
=\sum_{n\ge0}a_n t^n
=(\Gamma(\beta))^\gamma F^{(\gamma)}_{\alpha,\beta}(t).
\]
The right-hand side is finite for every \(t\ge0\). Therefore \(e^{|z|X}\) is integrable for every \(z\in\mathbb C\), and dominated convergence gives
\[
{\bf E}[e^{zX}]
=\sum_{n\ge0}\frac{z^n{\bf E}[X^n]}{n!}
=(\Gamma(\beta))^\gamma F^{(\gamma)}_{\alpha,\beta}(z),
\qquad z\in\mathbb C.
\]

This proves the claim.

For \(x>0\), substituting \(z=-x\) yields
\[
F^{(\gamma)}_{\alpha,\beta}(-x)
=\Gamma(\beta)^{-\gamma}{\bf E}[e^{-xX}].
\]
This is a positive Laplace transform. More explicitly, for every \(k\ge0\),
\[
(-1)^k\frac{d^k}{dx^k}F^{(\gamma)}_{\alpha,\beta}(-x)
=\Gamma(\beta)^{-\gamma}{\bf E}[X^k e^{-xX}]\ge0,
\]
where differentiating under the expectation is justified on compact subintervals of \((0,\infty)\) by the finite moment \({\bf E}[X^k]\).

This proves the claim.
\end{proof}

\begin{proposition}
\label{res:ab-reciprocal-sum-gate-a-plus-b-le-one-third}
If \(x\mapsto1/P_{a,b}(u,v;x)\) is completely monotone on \((0,\infty)\) for every \(v>u>0\), then \(a+b\le1/3\).
\end{proposition}
\begin{proof}
For \(a\ne b\), the Gini mean is
\[
G_{a,b}(u,v)
=\left(\frac{u^a+v^a}{u^b+v^b}\right)^{1/(a-b)}.
\]
The case \(a=b\) is obtained by continuity. Put \(v=u+h=u(1+t)\). Then
\[
\log \frac{G_{a,b}(u,u+h)}{u}
=\frac{1}{a-b}
\left[
\log(1+(1+t)^a)-\log(1+(1+t)^b)
\right].
\]
For a real parameter \(c\),
\[
\log(1+(1+t)^c)
=\log2+\frac c2t+\frac{c(c-2)}8t^2+O(t^3).
\]
Thus
\[
\log \frac{G_{a,b}(u,u+h)}{u}
=\frac12t+\frac{a+b-2}{8}t^2+O(t^3).
\]
Exponentiating gives
\[
G_{a,b}(u,u+h)
=u+\frac h2+\frac{a+b-1}{8u}h^2+O(h^3).
\]

Let \(Q_{a,b}=1/P_{a,b}\), \(r=a+b\), and \(y=x+u\). With \(v=u+h\),
\[
\log Q_{a,b}(u,u+h;x)
=\log\Gamma(y+h)-\log\Gamma(y)
-h\,\psi(x+G_{a,b}(u,u+h)).
\]
Using the Gini expansion,
\[
x+G_{a,b}(u,u+h)
=y+\frac h2+\frac{r-1}{8u}h^2+O(h^3).
\]
Taylor expansion gives
\[
\log\Gamma(y+h)-\log\Gamma(y)
=h\psi(y)+\frac{h^2}{2}\psi'(y)+\frac{h^3}{6}\psi''(y)+O(h^4),
\]
and
\[
h\psi\left(y+\frac h2+\frac{r-1}{8u}h^2+O(h^3)\right)
=h\psi(y)+\frac{h^2}{2}\psi'(y)
h^3\left(\frac{r-1}{8u}\psi'(y)+\frac18\psi''(y)\right)+O(h^4).
\]
Therefore
\[
\log Q_{a,b}(u,u+h;x)
=h^3\left(
\frac{\psi''(y)}{24}
-\frac{r-1}{8u}\psi'(y)
\right)+O(h^4).
\]

Assume \(x\mapsto Q_{a,b}(u,v;x)\) is completely monotone for every \(v>u>0\). Then it is decreasing and nonnegative. Standard Gamma and digamma asymptotics give
\[
\lim_{x\to\infty}Q_{a,b}(u,v;x)=1.
\]
Hence \(Q_{a,b}(u,v;x)\ge1\), so \(\log Q_{a,b}(u,v;x)\ge0\).

For fixed \(u,x>0\), let \(h\downarrow0\). The leading coefficient in the expansion must be nonnegative:
\[
\frac{\psi''(y)}{24}
-\frac{r-1}{8u}\psi'(y)
\ge0.
\]
Since \(\psi'(y)>0\),
\[
r\le 1+\frac{u\psi''(y)}{3\psi'(y)}.
\]
Here \(y=x+u\), so for each \(y>0\) and every \(0<u<y\),
\[
r\le 1+\frac{u\psi''(y)}{3\psi'(y)}.
\]
Letting \(u\uparrow y\) gives
\[
r\le 1+\frac{y\psi''(y)}{3\psi'(y)}
\qquad (y>0).
\]
Finally, as \(y\downarrow0\),
\[
\psi'(y)\sim y^{-2},
\qquad
\psi''(y)\sim -2y^{-3}.
\]
Therefore
\[
1+\frac{y\psi''(y)}{3\psi'(y)}\to \frac13,
\]
and hence
\[
a+b=r\le \frac13.
\]
\end{proof}

\begin{proposition}
\label{res:bz-gamma-quotient-derivative-normal-form}
For \(x>-1\), away from removable points, the sign of \(\widetilde F'(x)\) is the sign of \(N(x)=(x+6)(x^2+6)\psi(x+1)(\log(x^2+6)-\log(x+6))-(x^2+12x-6)\log\Gamma(x+1)\).
\end{proposition}
\begin{proof}
Set
\[
P(x)=(x+6)(x^2+6),\qquad Q(x)=x^2+12x-6.
\]
For \(x\ge8\), \(Q(x)>0\). The standard inequalities
\[
\psi(x+1)>\log x
\]
and the Robbins/Stirling upper bound imply
\[
\log\Gamma(x+1)\le x\log x
\qquad (x\ge8).
\]
Indeed the Robbins upper bound gives
\[
\log\Gamma(x+1)
<
\left(x+\frac12\right)\log x-x+\frac12\log(2\pi)+\frac1{12x},
\]
and the remaining excess
\[
\frac12\log x-x+\frac12\log(2\pi)+\frac1{12x}
\]
is already negative at \(x=8\) and decreasing thereafter.

It remains to lower-bound \(D(x)\). Let
\[
u(x)=\frac{x^2+6}{x+6}.
\]
For \(x\ge8\), \(u(x)>1\), so
\[
\log u(x)\ge \frac{2(u(x)-1)}{u(x)+1}.
\]
A direct rational simplification gives
\[
\frac{2(u(x)-1)}{u(x)+1}
-\frac{xQ(x)}{P(x)}
=
\frac{x^2(x^3-3x^2-18x-78)}
{(x^2+x+12)(x+6)(x^2+6)}.
\]
The denominator is positive, and \(x^3-3x^2-18x-78\) is increasing for \(x\ge8\) and has value \(98>0\) at \(x=8\). Hence
\[
D(x)=\log u(x)\ge \frac{xQ(x)}{P(x)}
\qquad (x\ge8).
\]

Combining the three inequalities,
\[
N(x)
=P(x)\psi(x+1)D(x)-Q(x)\log\Gamma(x+1)
>
P(x)(\log x)D(x)-Q(x)x\log x
\ge0.
\]
Therefore \(N(x)>0\), and consequently
\[
\widetilde F'(x)>0
\qquad (x\ge8).
\]
\end{proof}

The next result applies the preceding method to the fixed source problem \textbf{APP-0010}.
\begin{corollary}[APP-0010: Nielsen \(k\)-beta derivative-ratio monotonicity]
\label{res:nielsen-k-beta-derivative-ratio-monotonicity}
For \(k>0\), \(n\ge0\), and \(f_k(x)=x\beta_k(x)\), the ratio \(f_k^{(n+1)}(x)/(f_k^{(n)}(x)f_k^{(n+2)}(x))\) is strictly increasing on \((0,\infty)\) when \(n\) is odd and strictly decreasing on \((0,\infty)\) when \(n\) is even.

This resolves the fixed application label \textbf{APP-0010}: Nielsen \(k\)-beta derivative-ratio monotonicity
\emph{Source reference.} \cite{app-app-0010-source}.
\end{corollary}
\begin{proof}
Let \(\mu\) be a positive Borel measure on \([0,\infty)\) with finite moments after the tilt \(e^{-xt}\). Define
\[
a_j(x)=\int_0^\infty t^j e^{-xt}\,d\mu(t)>0
\]
for the needed indices. For fixed \(n\ge0\), set
\[
B_n(x)=\frac{a_{n+1}(x)}{a_n(x)a_{n+2}(x)}.
\]
Then \(B_n\) is strictly increasing whenever \(a_{n+1}(x)>0\).

Indeed \(a_j'(x)=-a_{j+1}(x)\). Put
\[
r_j(x)=\frac{a_{j+1}(x)}{a_j(x)}.
\]
Moment log-convexity, equivalently Cauchy--Schwarz applied to
\[
\int t^j e^{-xt}\,d\mu(t),
\]
gives \(r_{j+1}(x)\ge r_j(x)\). Differentiating,
\[
\frac{d}{dx}\log B_n(x)
=-\frac{a_{n+2}}{a_{n+1}}
+\frac{a_{n+1}}{a_n}
+\frac{a_{n+3}}{a_{n+2}}
=r_n(x)-r_{n+1}(x)+r_{n+2}(x).
\]
Since \(r_{n+2}(x)\ge r_{n+1}(x)\) and \(r_n(x)>0\), the logarithmic derivative is positive. Hence \(B_n\) is strictly increasing.

Let
\[
f_k(x)=x\beta_k(x).
\]
From the integral representation of \(\beta_k\), integration by parts gives
\[
f_k(x)
=x\int_0^\infty e^{-xt}\frac{dt}{1+e^{-kt}}
=\frac12+\int_0^\infty e^{-xt}\frac{k e^{-kt}}{(1+e^{-kt})^2}\,dt.
\]
Thus \(f_k\) is the Laplace transform of a positive measure consisting of an atom \(1/2\) at \(0\) plus the positive density
\[
w_k(t)=\frac{k e^{-kt}}{(1+e^{-kt})^2}
\]
on \((0,\infty)\). All tilted moments are finite, and \(a_{j}(x)>0\) for \(j\ge1\).

For \(a_j(x)=(-1)^j f_k^{(j)}(x)\), the bridge lemma gives that
\[
B_{k,n}(x)=\frac{a_{n+1}(x)}{a_n(x)a_{n+2}(x)}
\]
is strictly increasing for every \(k>0\), \(n\ge0\), and \(x>0\). Since
\[
R_{k,n}(x)=(-1)^{n+1}B_{k,n}(x),
\]
we get the precise monotonicity law:

if \(n\) is odd, \(R_{k,n}\) is strictly increasing on \((0,\infty)\);
if \(n\) is even, \(R_{k,n}\) is strictly decreasing on \((0,\infty)\).

This solves the Yin--Zhang Nielsen \(k\)-beta derivative-ratio open problem in the parity-refined form.
\end{proof}

\begin{proposition}[Complete monotonicity of \(V_q\)]
\label{res:baricz-vq-q-complete-monotonicity}
For every fixed \(x>0\), the function \(a\mapsto V_{a-1}(x)\) is completely monotone on \((0,\infty)\).
\end{proposition}
\begin{proof}
Put \(u=t^2-x^2\). Then \(t=(u+x^2)^{1/2}\) and
\[
dt=\frac{du}{2(u+x^2)^{1/2}}.
\]
Since \(e^{x^2}e^{-t^2}=e^{-u}\), the factor \(2\) in the definition cancels the \(1/2\) from \(dt\), giving
\[
V_q(x)=\frac1{\Gamma(q+1)}
\int_0^\infty e^{-u}u^q(u+x^2)^{-1/2}\,du.
\]
This integral is finite for \(q>-1\): near \(0\) the integrand is \(O(u^q)\), and at infinity it is exponentially decaying.

For \(r>0\),
\[
r^{-1/2}=\frac1{\sqrt\pi}\int_0^\infty s^{-1/2}e^{-rs}\,ds.
\]
With \(r=u+x^2\), Tonelli's theorem applies because the integrand is nonnegative. Thus
\[
\begin{aligned}
V_q(x)
&=\frac1{\Gamma(q+1)\sqrt\pi}
\int_0^\infty\int_0^\infty
e^{-u}u^q s^{-1/2}e^{-(u+x^2)s}\,ds\,du\\
&=\frac1{\Gamma(q+1)\sqrt\pi}
\int_0^\infty s^{-1/2}e^{-x^2s}
\left(\int_0^\infty u^q e^{-(1+s)u}\,du\right)\,ds\\
&=\frac1{\sqrt\pi}\int_0^\infty
s^{-1/2}e^{-x^2s}(1+s)^{-(q+1)}\,ds.
\end{aligned}
\]
The last step uses
\[
\int_0^\infty u^q e^{-(1+s)u}\,du
=\frac{\Gamma(q+1)}{(1+s)^{q+1}}.
\]

This proves the claim.

Let \(a=q+1>0\). The normal form becomes
\[
V_{a-1}(x)=\frac1{\sqrt\pi}\int_0^\infty
s^{-1/2}e^{-x^2s}(1+s)^{-a}\,ds.
\]
For each integer \(n\ge0\), differentiation under the integral gives
\[
(-1)^n\frac{d^n}{da^n}V_{a-1}(x)
=\frac1{\sqrt\pi}\int_0^\infty
(\log(1+s))^n s^{-1/2}e^{-x^2s}(1+s)^{-a}\,ds\ge0.
\]
The differentiation is justified locally uniformly in \(a>0\): near \(s=0\), \((\log(1+s))^n s^{-1/2}=O(s^{n-1/2})\), and as \(s\to\infty\), the factor \(e^{-x^2s}\) dominates every logarithmic and polynomial factor.

Thus \(a\mapsto V_{a-1}(x)\) is completely monotone on \((0,\infty)\). Equivalently, after the change of variables \(y=\log(1+s)\), it is a positive Laplace transform in \(a\).

This proves the claim.

For \(a>0\), write
\[
F_x(a)=V_{a-1}(x)
=\int_{(0,\infty)} e^{-ay}\,d\nu_x(y),
\]
where \(\nu_x\) is the positive pushforward of
\[
\frac1{\sqrt\pi}s^{-1/2}e^{-x^2s}\,ds
\]
under \(y=\log(1+s)\). This measure is not concentrated at one point because it has positive density on \(s>0\), hence on \(y>0\).

For \(a_1\ne a_2\) and \(0<\lambda<1\), strict Holder inequality gives
\[
F_x(\lambda a_1+(1-\lambda)a_2)
<
F_x(a_1)^\lambda F_x(a_2)^{1-\lambda}.
\]
Strictness holds because \(e^{-a_1y}\) and \(e^{-a_2y}\) are not proportional \(\nu_x\)-almost everywhere unless \(\nu_x\) is supported at a single \(y\).

Thus \(a\mapsto F_x(a)\) is strictly log-convex on \((0,\infty)\). Since \(a=q+1\), the shift preserves strict log-convexity, so \(q\mapsto V_q(x)\) is strictly log-convex on \((-1,\infty)\) for every fixed \(x>0\).

This proves the claim.
\end{proof}

The next result applies the preceding method to the fixed source problem \textbf{APP-0022}.
\begin{corollary}[APP-0022: Chen-Choi Gurland gamma-ratio conjecture is logarithmically completely monotone]
\label{res:chen-choi-gurland-logf-strict-cm}
For \(F(x)=\Gamma(1/x)\Gamma(3/x)/\Gamma(2/x)^2\), the function \(x\mapsto\log F(x)\) is strictly completely monotone on \((0,\infty)\).

This resolves the fixed application label \textbf{APP-0022}: Prove logarithmic complete monotonicity of the Gurland gamma ratio conjectured in the source.
\emph{Source reference.} \cite{app-app-0022-source}.
\end{corollary}
\begin{proof}
Set
\[
  A(u)=2\log(u+2)-\log(u+1)-\log(u+3).
\]
The identity
\[
 A(u)=\int_0^\infty e^{-ut}{e^{-t}(1-e^{-t})^2\over t}\,dt
\]
shows that \(A\) is strictly completely monotone.  The Weierstrass product for
\(\Gamma\) gives
\[
  \log F(x)=\log {4\over3}+\sum_{m\ge1}A(mx).
\]
Termwise differentiation is justified on compact subintervals of
\((0,\infty)\), and each differentiated summand has the required strict
alternating sign.  This proves logarithmic complete monotonicity.
\end{proof}

\begin{proposition}
\label{res:prabhakar-q-stable-subordination-normal-form}
For \(0<\alpha<1\), \(\gamma>0\), \(\beta>\alpha\gamma\), and \(\lambda>0\), let \(P^\gamma_{\alpha,\beta}\) be Sibisi's transform-normalized Pollard measure satisfying \(E^\gamma_{\alpha,\beta}(-u)=\int_0^\infty e^{-ur}\,dP^\gamma_{\alpha,\beta}(r)\). Then the measure \(Q^\gamma_{\alpha,\beta}(A\mid\lambda)=\int\Pr((\lambda r)^{1/\alpha}S_\alpha\in A)\,dP^\gamma_{\alpha,\beta}(r)\) satisfies \(E^\gamma_{\alpha,\beta}(-\lambda x^\alpha)=\int_0^\infty e^{-xt}\,dQ^\gamma_{\alpha,\beta}(t\mid\lambda)\).
\end{proposition}
\begin{proof}
Let \(P\) be a finite positive measure on \([0,\infty)\), let
\[
F(u)=\int_0^\infty e^{-ur}\,dP(r),
\]
let \(0<\alpha<1\), \(\lambda>0\), and let \(S_\alpha\) be positive \(\alpha\)-stable:
\[
\mathbb E e^{-xS_\alpha}=e^{-x^\alpha}.
\]
Define a finite positive measure \(Q\) by
\[
Q(A)=\int_{[0,\infty)}
\Pr\{(\lambda r)^{1/\alpha}S_\alpha\in A\}\,dP(r).
\]
Then Tonelli's theorem gives, for \(x>0\),
\[
\int_0^\infty e^{-xt}\,dQ(t)
=\int_{[0,\infty)}
\mathbb E e^{-x(\lambda r)^{1/\alpha}S_\alpha}\,dP(r)
=\int_{[0,\infty)}e^{-\lambda r x^\alpha}\,dP(r)
=F(\lambda x^\alpha).
\]

Apply the lemma to Sibisi's transform-normalized \(P=P^\gamma_{\alpha,\beta}\) and
\[
F(u)=E^\gamma_{\alpha,\beta}(-u).
\]
For \(0<\alpha<1\), \(\gamma>0\), \(\beta>\alpha\gamma\), and \(\lambda>0\), set
\[
Q^\gamma_{\alpha,\beta}(A\mid\lambda)
=\int_{[0,\infty)}
\Pr\{(\lambda r)^{1/\alpha}S_\alpha\in A\}\,
dP^\gamma_{\alpha,\beta}(r).
\]
Then
\[
\int_0^\infty e^{-xt}\,dQ^\gamma_{\alpha,\beta}(t\mid\lambda)
=E^\gamma_{\alpha,\beta}(-\lambda x^\alpha).
\]
By uniqueness of finite Laplace transforms, this is Sibisi's \(Q^\gamma_{\alpha,\beta}(\cdot\mid\lambda)\).

If \(S_\alpha\) has density \(f_\alpha\), then the non-atomic density part is
\[
q^\gamma_{\alpha,\beta}(t\mid\lambda)
=\int_{(0,\infty)}
(\lambda r)^{-1/\alpha}
f_\alpha\!\left(\frac{t}{(\lambda r)^{1/\alpha}}\right)
dP^\gamma_{\alpha,\beta}(r),
\]
with an atom \(P^\gamma_{\alpha,\beta}(\{0\})\delta_0\) if \(P^\gamma_{\alpha,\beta}\) has an atom at zero. When \(P^\gamma_{\alpha,\beta}\) is represented by Sibisi's Pollard density \(p^\gamma_{\alpha,\beta}(r)\), this becomes
\[
q^\gamma_{\alpha,\beta}(t\mid\lambda)
=\int_0^\infty
(\lambda r)^{-1/\alpha}
f_\alpha\!\left(\frac{t}{(\lambda r)^{1/\alpha}}\right)
p^\gamma_{\alpha,\beta}(r)\,dr.
\]

Because
\[
E^\gamma_{\alpha,\beta}(0)=\frac1{\Gamma(\beta)},
\]
the representing measures \(P^\gamma_{\alpha,\beta}\) and \(Q^\gamma_{\alpha,\beta}(\cdot\mid\lambda)\) have total mass \(1/\Gamma(\beta)\). They are finite positive measures, not generally probability measures. The normalized probability version is
\[
\widehat Q^\gamma_{\alpha,\beta}=\Gamma(\beta)Q^\gamma_{\alpha,\beta},
\qquad
\mathcal L\widehat Q^\gamma_{\alpha,\beta}(x)
=\Gamma(\beta)E^\gamma_{\alpha,\beta}(-\lambda x^\alpha).
\]

The boundary \(\beta=\alpha\gamma\) is not included here because Sibisi's ordinary convolution density uses \(1/\Gamma(\beta-\alpha\gamma)\). A limiting treatment is a separate open frontier.
\end{proof}

The next result applies the preceding method to the fixed source problem \textbf{APP-0012}.
\begin{corollary}[APP-0012: Baricz \(\(V_q\)\) strict \(q\)-log-convexity]
\label{res:baricz-vq-strict-q-logconvexity}
For every fixed \(x>0\), the function \(q\mapsto V_q(x)\) is strictly log-convex on \((-1,\infty)\).

This resolves the fixed application label \textbf{APP-0012}: Baricz \(V\_q\) strict \(q\)-log-convexity
\emph{Source reference.} \cite{app-app-0012-source}.
\end{corollary}
\begin{proof}
Put \(u=t^2-x^2\).  Then
\begin{equation}
\label{eq:baricz-u-normal-form}
V_q(x)=\frac1{\Gamma(q+1)}
\int_0^\infty e^{-u}u^q(u+x^2)^{-1/2}\,du.
\end{equation}
Using
\[
r^{-1/2}=\frac1{\sqrt\pi}\int_0^\infty s^{-1/2}e^{-rs}\,ds
\]
and Tonelli's theorem, we obtain
\begin{equation}
\label{eq:baricz-laplace-normal-form}
V_q(x)=\frac1{\sqrt\pi}\int_0^\infty
s^{-1/2}e^{-x^2s}(1+s)^{-(q+1)}\,ds.
\end{equation}
Let \(a=q+1>0\).  Differentiation under the integral gives
\begin{equation}
\label{eq:baricz-complete-monotonicity}
(-1)^n\frac{d^n}{da^n}V_{a-1}(x)
=\frac1{\sqrt\pi}\int_0^\infty
(\log(1+s))^n s^{-1/2}e^{-x^2s}(1+s)^{-a}\,ds\ge0.
\end{equation}
Equivalently, after the change of variables \(y=\log(1+s)\), \(V_{a-1}(x)\) is a positive Laplace transform in \(a\).  The representing measure is not concentrated at a single point, because the density in \eqref{eq:baricz-laplace-normal-form} is positive on \(s>0\).  Strict Holder inequality for two different exponents therefore gives strict log-convexity in \(a\), and the shift \(a=q+1\) gives the stated strict log-convexity in \(q\).
\end{proof}

\begin{proposition}
\label{res:baricz-vq-q-logconvexity-open-problem}
For Baricz's function \(V_q(x)=\frac{2e^{x^2}}{\Gamma(q+1)}\int_x^\infty e^{-t^2}(t^2-x^2)^q\,dt\), prove that for every fixed \(x>0\), the map \(q\mapsto V_q(x)\) is log-convex on \((-1,\infty)\).
\end{proposition}
\begin{proof}
Put \(u=t^2-x^2\). Then \(t=(u+x^2)^{1/2}\) and
\[
dt=\frac{du}{2(u+x^2)^{1/2}}.
\]
Since \(e^{x^2}e^{-t^2}=e^{-u}\), the factor \(2\) in the definition cancels the \(1/2\) from \(dt\), giving
\[
V_q(x)=\frac1{\Gamma(q+1)}
\int_0^\infty e^{-u}u^q(u+x^2)^{-1/2}\,du.
\]
This integral is finite for \(q>-1\): near \(0\) the integrand is \(O(u^q)\), and at infinity it is exponentially decaying.

For \(r>0\),
\[
r^{-1/2}=\frac1{\sqrt\pi}\int_0^\infty s^{-1/2}e^{-rs}\,ds.
\]
With \(r=u+x^2\), Tonelli's theorem applies because the integrand is nonnegative. Thus
\[
\begin{aligned}
V_q(x)
&=\frac1{\Gamma(q+1)\sqrt\pi}
\int_0^\infty\int_0^\infty
e^{-u}u^q s^{-1/2}e^{-(u+x^2)s}\,ds\,du\\
&=\frac1{\Gamma(q+1)\sqrt\pi}
\int_0^\infty s^{-1/2}e^{-x^2s}
\left(\int_0^\infty u^q e^{-(1+s)u}\,du\right)\,ds\\
&=\frac1{\sqrt\pi}\int_0^\infty
s^{-1/2}e^{-x^2s}(1+s)^{-(q+1)}\,ds.
\end{aligned}
\]
The last step uses
\[
\int_0^\infty u^q e^{-(1+s)u}\,du
=\frac{\Gamma(q+1)}{(1+s)^{q+1}}.
\]

This proves the claim.

Let \(a=q+1>0\). The normal form becomes
\[
V_{a-1}(x)=\frac1{\sqrt\pi}\int_0^\infty
s^{-1/2}e^{-x^2s}(1+s)^{-a}\,ds.
\]
For each integer \(n\ge0\), differentiation under the integral gives
\[
(-1)^n\frac{d^n}{da^n}V_{a-1}(x)
=\frac1{\sqrt\pi}\int_0^\infty
(\log(1+s))^n s^{-1/2}e^{-x^2s}(1+s)^{-a}\,ds\ge0.
\]
The differentiation is justified locally uniformly in \(a>0\): near \(s=0\), \((\log(1+s))^n s^{-1/2}=O(s^{n-1/2})\), and as \(s\to\infty\), the factor \(e^{-x^2s}\) dominates every logarithmic and polynomial factor.

Thus \(a\mapsto V_{a-1}(x)\) is completely monotone on \((0,\infty)\). Equivalently, after the change of variables \(y=\log(1+s)\), it is a positive Laplace transform in \(a\).

This proves the claim.

For \(a>0\), write
\[
F_x(a)=V_{a-1}(x)
=\int_{(0,\infty)} e^{-ay}\,d\nu_x(y),
\]
where \(\nu_x\) is the positive pushforward of
\[
\frac1{\sqrt\pi}s^{-1/2}e^{-x^2s}\,ds
\]
under \(y=\log(1+s)\). This measure is not concentrated at one point because it has positive density on \(s>0\), hence on \(y>0\).

For \(a_1\ne a_2\) and \(0<\lambda<1\), strict Holder inequality gives
\[
F_x(\lambda a_1+(1-\lambda)a_2)
<
F_x(a_1)^\lambda F_x(a_2)^{1-\lambda}.
\]
Strictness holds because \(e^{-a_1y}\) and \(e^{-a_2y}\) are not proportional \(\nu_x\)-almost everywhere unless \(\nu_x\) is supported at a single \(y\).

Thus \(a\mapsto F_x(a)\) is strictly log-convex on \((0,\infty)\). Since \(a=q+1\), the shift preserves strict log-convexity, so \(q\mapsto V_q(x)\) is strictly log-convex on \((-1,\infty)\) for every fixed \(x>0\).

This proves the claim.
\end{proof}

\begin{theorem}
\label{res:simon-gamma-quotient-derivative-positive-laplace-density}
For every \(0<\alpha<1\), the derivative of \(F_\alpha(x)=\Gamma(x+\alpha)/(\Gamma(x)x^\alpha)\) is completely monotone, with a positive Laplace-density representation obtained from the decreasing kernel \(J_\alpha\).
\end{theorem}
\begin{proof}
For \(0<\alpha<1\), put
\[
F_\alpha(x)=\frac{\Gamma(x+\alpha)}{\Gamma(x)x^\alpha},\qquad x>0.
\]
Then \(F_\alpha\) is a Bernstein function. In fact, \(1-F_\alpha\) is completely monotone.

The stronger route asserting that \(F_\alpha\) is a complete Bernstein function is false.

Let \(\beta=1-\alpha\in(0,1)\). Since \(\Gamma(x+1)=x\Gamma(x)\), the beta integral gives
\[
F_\alpha(x)
=x^\beta\frac{\Gamma(x+\alpha)}{\Gamma(x+1)}
=\frac{x^\beta}{\Gamma(\beta)}
\int_0^\infty e^{-xu}(e^u-1)^{-\alpha}\,du.
\]
Also
\[
1=\frac{x^\beta}{\Gamma(\beta)}
\int_0^\infty e^{-xu}u^{-\alpha}\,du.
\]
Thus
\[
1-F_\alpha(x)
=\frac{x^\beta}{\Gamma(\beta)}
\int_0^\infty e^{-xu}\left(u^{-\alpha}-(e^u-1)^{-\alpha}\right)\,du.
\]

Define
\[
J_\alpha(t)=\int_0^t (t-u)^{\alpha-1}(e^u-1)^{-\alpha}\,du.
\]
By Laplace convolution,
\[
\mathcal L[J_\alpha](x)
=\Gamma(\alpha)x^{-\alpha}
\int_0^\infty e^{-xu}(e^u-1)^{-\alpha}\,du.
\]
Therefore
\[
F_\alpha(x)
=\frac{x}{\Gamma(\alpha)\Gamma(1-\alpha)}\mathcal L[J_\alpha](x).
\]

After the substitution \(u=tv\),
\[
J_\alpha(t)=
\int_0^1
(1-v)^{\alpha-1}
\left(\frac{t}{e^{tv}-1}\right)^\alpha\,dv.
\]
For fixed \(v\in(0,1]\), the function
\[
t\mapsto \frac{t}{e^{tv}-1}
\]
is decreasing because
\[
\frac{d}{dt}\frac{t}{e^{tv}-1}
=
\frac{e^{tv}-1-tv e^{tv}}{(e^{tv}-1)^2}\le0.
\]
The integrand is dominated by \((1-v)^{\alpha-1}v^{-\alpha}\), which is integrable on \((0,1)\). Hence \(J_\alpha\) is decreasing. Moreover,
\[
J_\alpha(0+)
=\int_0^1(1-v)^{\alpha-1}v^{-\alpha}\,dv
=\Gamma(\alpha)\Gamma(1-\alpha),
\]
and \(J_\alpha(t)\to0\) as \(t\to\infty\).

Let \(\mu_\alpha=-dJ_\alpha\). This is a positive finite measure with total mass
\[
\mu_\alpha([0,\infty))=\Gamma(\alpha)\Gamma(1-\alpha).
\]
Stieltjes integration by parts gives, for \(x>0\),
\[
\int_0^\infty e^{-xt}\,\mu_\alpha(dt)
=
\Gamma(\alpha)\Gamma(1-\alpha)-x\mathcal L[J_\alpha](x).
\]
Using the formula for \(F_\alpha\),
\[
1-F_\alpha(x)
=
\frac{1}{\Gamma(\alpha)\Gamma(1-\alpha)}
\int_0^\infty e^{-xt}\,\mu_\alpha(dt).
\]
Thus \(1-F_\alpha\) is completely monotone.

Consequently,
\[
F_\alpha'(x)
=
\frac{1}{\Gamma(\alpha)\Gamma(1-\alpha)}
\int_0^\infty t e^{-xt}\,\mu_\alpha(dt),
\]
so \(F_\alpha'\) is completely monotone. Since \(F_\alpha(x)>0\), \(F_\alpha\) is a Bernstein function.

For \(t>0\), differentiating the \(v\)-integral under the same endpoint domination gives \(\mu_\alpha(dt)=m_\alpha(t)\,dt\), where
\[
m_\alpha(t)
=
\alpha\int_0^1
(1-v)^{\alpha-1}
\left(\frac{t}{e^{tv}-1}\right)^\alpha
\left(
\frac{v e^{tv}}{e^{tv}-1}-\frac1t
\right)\,dv.
\]
The bracket is positive, because with \(a=tv>0\),
\[
\frac{a e^a}{e^a-1}>1.
\]
Therefore the Laplace-density route in the argument is also closed.

The case \(\alpha=1/2\) is the specialization
\[
F_{1/2}(x)=\frac{\Gamma(x+1/2)}{\Gamma(x)\sqrt{x}},
\]
and the same complement-CM representation proves that \(F_{1/2}\) is Bernstein.

For \(z=-r+i0\) with \(r\in(\alpha,1)\), take the upper-lip branch of \(z^{-\alpha}\). Then
\[
F_\alpha(-r+i0)
=
\frac{\Gamma(\alpha-r)}{\Gamma(-r)}r^{-\alpha}e^{-i\pi\alpha}.
\]
Here \(\Gamma(\alpha-r)<0\) and \(\Gamma(-r)<0\), so the real gamma ratio is positive. Hence
\[
\operatorname{Im} F_\alpha(-r+i0)<0.
\]
A complete Bernstein function has the Pick property and maps the upper half-plane into itself. Therefore \(F_\alpha\) is not a complete Bernstein function for \(0<\alpha<1\).
\end{proof}

\begin{proposition}
\label{res:simon-gamma-quotient-bf-alpha-window-open-problem}
For every \(0<\alpha<1\), the function \(F_\alpha(x)=\Gamma(x+\alpha)/(\Gamma(x)x^\alpha)\) is a Bernstein function on \((0,\infty)\).
\end{proposition}
\begin{proof}
Source artifacts:

Gomilko--Tomilov source: arXiv:2210.03474, "To the theory of generalized Stieltjes transforms"
Simon source: Thomas Simon, "Moment problems related to Bernstein functions", Annales de la Faculte des sciences de Toulouse 29 (2020), 577--594, DOI 10.5802/afst.1640

For \(\alpha,\beta>0\), \(s,t\ge0\), and \(x>0\), the beta identity gives
\[
(x+s)^{-\alpha}(x+t)^{-\beta}
=
\frac{\Gamma(\alpha+\beta)}{\Gamma(\alpha)\Gamma(\beta)}
\int_0^1
\frac{u^{\alpha-1}(1-u)^{\beta-1}}
{\bigl(x+us+(1-u)t\bigr)^{\alpha+\beta}}
\,du.
\]

If
\[
f(x)=\int_{[0,\infty)}(x+s)^{-\alpha}\,d\mu(s),
\qquad
g(x)=\int_{[0,\infty)}(x+t)^{-\beta}\,d\nu(t),
\]
then Tonelli gives
\[
f(x)g(x)=\int_{[0,\infty)}(x+r)^{-(\alpha+\beta)}\,d\kappa(r),
\]
where \(\kappa\) is the positive pushforward measure defined by
\[
\int \varphi(r)\,d\kappa(r)
=
\frac{\Gamma(\alpha+\beta)}{\Gamma(\alpha)\Gamma(\beta)}
\int\!\!\int\!\!\int_0^1
\varphi\bigl(us+(1-u)t\bigr)
u^{\alpha-1}(1-u)^{\beta-1}
\,du\,d\mu(s)\,d\nu(t).
\]

Thus pure generalized Stieltjes transforms satisfy
\[
\mathcal S_\alpha^{0}\mathcal S_\beta^{0}\subseteq \mathcal S_{\alpha+\beta}^{0}.
\]
With the usual nonnegative constant term convention, constants and lower-order terms require the standard order-lift convention already recorded in the source vocabulary; no application should rely on constants being absorbed unless that convention is explicitly cited.

For \(0<\alpha<1\), set
\[
F_\alpha(x)=\frac{\Gamma(x+\alpha)}{\Gamma(x)x^\alpha},
\qquad
L_\alpha(x)=\log F_\alpha(x).
\]
Then
\[
L_\alpha'(x)
=
\psi(x+\alpha)-\psi(x)-\frac{\alpha}{x}
=
\int_0^\infty e^{-xt}
\left(
\frac{1-e^{-\alpha t}}{1-e^{-t}}-\alpha
\right)\,dt.
\]
The kernel is positive for \(t>0\), because \(u\mapsto 1-e^{-ut}\) is strictly concave on \([0,1]\), hence
\[
1-e^{-\alpha t}>\alpha(1-e^{-t}).
\]
Therefore \(L_\alpha'\) is completely monotone, and \(1/F_\alpha\) is logarithmically completely monotone.

This does not prove that \(F_\alpha\) is a Bernstein function. Simon's source explicitly states that the Bernstein character of the corresponding \(\Phi(\lambda)=\Gamma(1-t+\lambda)/(\lambda^{1-t}\Gamma(\lambda))\) is a puzzling question, and also says the reciprocal LCM property is necessary but not sufficient.
\end{proof}
